
\section{Subgroups of type \textup{(C)} of group preschemes over an
arbitrary prescheme}
\SGARefTarget{sga3:xiv:section:3}

\paragraph{\textbf{Theorem 3.1.}}
\SGARefTarget{sga3:xiv:3.1}
{\itshape Let \(S\) be a prescheme, let \(G\) be a smooth
\(S\)-group prescheme, and let \(\mathfrak g\) be its Lie algebra,
which is a locally free module of finite type over \(S\). Let
\(\mathfrak h\) be a Lie subalgebra of \(\mathfrak g\) that, as a
module, is locally a direct factor in \(\mathfrak g\), and suppose
that, for every \(s\in S\), the geometric fiber
\(\mathfrak h_{\bar s}\) contains a Cartan subalgebra of
\(\mathfrak g_{\bar s}\). Let \(a\) be a quasi-regular section of
\(\mathfrak g\) \textup{(\SGARefLink{sga3:xiv:2.5}{2.5})}. Then
\[
  M_a=\operatorname{Transp}_G(a,\mathfrak h),
\]
the subfunctor of \(G\) whose \(S'\)-valued points are the
\(g\in G(S')\) such that
\[
  \operatorname{ad}(g)\cdot a_{S'}
    \in\Gamma(S',\mathfrak h_{S'}),
\]
is represented by a closed subprescheme of \(G\), smooth over
\(S\), whose structure morphism to \(S\) is surjective.}
\footnote{Editors' note: This transporter is denoted
\(\operatorname{Transp}_G(a,\mathfrak h)\).}

\paragraph{\textit{Proof.}}
Consider the canonical morphism
\[
  \varphi\colon G\times_S W(\mathfrak h)\longrightarrow
  W(\mathfrak g),
  \qquad
  (g,x)\longmapsto\operatorname{ad}(g)\cdot x.
\]
Then \(M_a\) is \(S\)-isomorphic to \(\varphi^{-1}(a)\), the inverse
image under \(\varphi\) of \(a\), viewed as a section of
\(W(\mathfrak g)\) over \(S\). To prove that \(M_a\) is smooth, it
is enough to show that \(\varphi\) is smooth at the points of
\(G\times_S W(\mathfrak h)\) lying over \(\operatorname{Im}(a)\).
More generally, \(\varphi\) is smooth at every point lying over a
point of \(W(\mathfrak g)\) that is regular in its fiber
\(W(\mathfrak g(s))\) over \(S\).

Indeed, the source and target of \(\varphi\) are smooth over \(S\),
and hence are flat and locally of finite presentation over \(S\).
We may therefore check the assertion fiber by fiber. This reduces
us to the case in which \(S\) is the spectrum of an algebraically
closed field \(k\), \(G\) is a locally algebraic group over \(k\),
\(\mathfrak h\) is a Lie subalgebra of its Lie algebra
\(\mathfrak g\) containing a Cartan subalgebra of \(\mathfrak g\),
and \(a\) is a regular point of \(\mathfrak g\).

Since \(\varphi\) is a \(G\)-morphism, we may suppose that the point
of \(G\times W(\mathfrak h)\) under consideration is of the form
\((e,a)\). We may also suppose that \(G\) is connected, and hence
of finite type over \(k\). The assertion is then precisely
\SGARefLink{sga3:xiii:corollary:5:4}{XIII~5.4}.

It is also immediate that \(M_a\) is a closed subprescheme of \(G\)
of finite presentation over \(S\): it is the inverse image of
\(W(\mathfrak h)\) under the morphism
\[
  G\longrightarrow W(\mathfrak g),
  \qquad
  g\longmapsto\operatorname{ad}(g)\cdot a.
\]
Finally, the surjectivity of \(M_a\to S\) likewise reduces to the
case of an algebraically closed base field. By hypothesis,
\(\mathfrak h\) then contains a Cartan subalgebra
\(\mathfrak d'\). By the conjugacy theorem \SGARefLink{sga3:xiii:theorem:6:1}{XIII~6.1}\textup{(a)},
\(\mathfrak d'\) is conjugate to
\[
  \mathfrak d=\Nil(a,\mathfrak g).
\]
It follows that \(\mathfrak h\) contains a conjugate of \(a\).
This completes the proof.
